% Tiny little diagrams for integral subscripts

\newcommand{\scriptdiagrscale}{0.3}
\newlength{\scriptdiagrleglen}\setlength{\scriptdiagrleglen}{0.2cm}
\newlength{\scriptdiagrleg}\setlength{\scriptdiagrleg}{\scriptdiagrleglen}\addtolength{\scriptdiagrleg}{\bubblesize}
\newcommand{\tad}{%
    {\tikzset{external/export next=false}\tikz[scale=\scriptdiagrscale, transform shape]{%
        \tikzset{pion/.style={ thin }}
        \path[use as bounding box] (-1.5\scriptdiagrleglen,0) rectangle (+1.5\scriptdiagrleglen,.6\tadpolesize);
        \draw[pion] (-1.5\scriptdiagrleglen,0) -- (0,0) pic {tadpole} -- (+1.5\scriptdiagrleglen,0);}}}
\newcommand{\bub}{%
    {\tikzset{external/export next=false}\tikz[scale=\scriptdiagrscale, transform shape]{%
        \tikzset{pion/.style={ thin }}
        \draw (0,0) pic {bubble};
        % This (0,0) +  avoids coordinate shenanigans in tikzpicutres
        \draw[pion] (0,0) + (-\bubblesize,0) -- +(-\scriptdiagrleg,0);
        \draw[pion] (0,0) + (+\bubblesize,0) -- +(+\scriptdiagrleg,0);
        }}}
\newcommand{\sunset}{%
    {\tikzset{external/export next=false}\tikz[scale=\scriptdiagrscale, transform shape]{%
        \tikzset{pion/.style={ thin }}
        \renewcommand{\bubbleaspect}{1}
        \draw[pion] (0,0) -- (-\bubblesize,0) -- +(-\scriptdiagrleglen,0);
        \draw[pion] (0,0) -- (+\bubblesize,0) -- +(+\scriptdiagrleglen,0);
        \draw (0,0) pic {bubble};
        }}}
\newcommand{\bball}{%
    {\tikzset{external/export next=false}\tikz[scale=\scriptdiagrscale, transform shape]{%
        \tikzset{pion/.style={ thin }}
        \renewcommand{\bubbleaspect}{1}
        \draw[pion] (0,0) + (-\bubblesize,0) -- +(-\scriptdiagrleg,0);
        \draw[pion] (0,0) + (+\bubblesize,0) -- +(+\scriptdiagrleg,0);
        \draw (0,0) pic {bubble} pic[yscale=.4] {bubble};
        }}}
\newcommand{\fatcat}{%
    {\tikzset{external/export next=false}\tikz[scale=\scriptdiagrscale, transform shape]{%
        \tikzset{pion/.style={ thin }}
        \renewcommand{\bubbleaspect}{1}
        \draw[pion] (0,0) + (-\bubblesize,0) -- +(-\scriptdiagrleg,0);
        \draw[pion] (0,0) + (+\bubblesize,0) -- +(+\scriptdiagrleg,0);
        \draw (0,0) pic {bubble=40};
        \draw (node40) pic[rotate=-40, scale=.6] {bubble};
        }}}
\newcommand{\catseye}{%
    {\tikzset{external/export next=false}\tikz[scale=\scriptdiagrscale, transform shape]{%
        \tikzset{pion/.style={ thin }}
        \renewcommand{\bubbleaspect}{1}
        \draw[pion] (0,0) + (-\bubblesize,0) -- +(-\scriptdiagrleg,0);
        \draw[pion] (0,0) + (+\bubblesize,0) -- +(+\scriptdiagrleg,0);
        \draw (0,0) pic {bubble} pic[xscale=.4] {bubble};
        }}}

\newcommand{\catseyeplus}{%
    {\tikzset{external/export next=false}\tikz[scale=\scriptdiagrscale, transform shape]{%
        \tikzset{pion/.style={ thin }}
        \renewcommand{\bubbleaspect}{1}
        \draw[pion] (0,0) + (-\bubblesize,0) -- +(-\scriptdiagrleg,0);
        \draw[pion] (0,0) + (+\bubblesize,0) -- +(+\scriptdiagrleg,0);
        \draw (0,0) pic {bubble={60,120,270}};
        \draw[pion] (node60)  -- (node270);
        \draw[pion] (node120) -- (node270);
        }}}
\newcommand{\catseyeplusplus}{%
    {\tikzset{external/export next=false}\tikz[scale=\scriptdiagrscale, transform shape]{%
        \tikzset{pion/.style={ thin }}
        \renewcommand{\bubbleaspect}{1}
        \draw[pion] (0,0) + (-\bubblesize,0) -- +(-\scriptdiagrleg,0);
        \draw[pion] (0,0) + (+\bubblesize,0) -- +(+\scriptdiagrleg,0);
        \draw (0,0) pic {bubble={60,120,240,300}};
        \draw[pion] (node60)  to[bend right=30] (node300);
        \draw[pion] (node240) to[bend right=30] (node120);
        }}}
